%%
%% Copyright 2021 OXFORD UNIVERSITY PRESS
%%
%% This file is part of the 'ima-authoring-template Bundle'.
%% ---------------------------------------------
%%
%% It may be distributed under the conditions of the LaTeX Project Public
%% License, either version 1.2 of this license or (at your option) any
%% later version.  The latest version of this license is in
%%    http://www.latex-project.org/lppl.txt
%% and version 1.2 or later is part of all distributions of LaTeX
%% version 1999/12/01 or later.
%%
%% The list of all files belonging to the 'ima-authoring-template Bundle' is
%% given in the file `manifest.txt'.
%%
%% Template article for OXFORD UNIVERSITY PRESS's document class `ima-authoring-template'
%% with bibliographic references
%%

\documentclass[numbers,webpdf,imaiai]{ima-authoring-template}%
%\documentclass[namedate,webpdf,imaman]{ima-authoring-template}

%\usepackage{showframe}

\graphicspath{{Fig/}}

% line numbers
%\usepackage[mathlines, switch]{lineno}
%\usepackage[right]{lineno}

\theoremstyle{thmstyletwo}%
\newtheorem{theorem}{Theorem}%  meant for continuous numbers
%%\newtheorem{theorem}{Theorem}[section]% meant for sectionwise numbers
%% optional argument [theorem] produces theorem numbering sequence instead of independent numbers for Proposition
\newtheorem{proposition}[theorem]{Proposition}%
%%\newtheorem{proposition}{Proposition}% to get separate numbers for theorem and proposition etc.

\newtheorem{example}{Example}%
\newtheorem{remark}{Remark}%

\newtheorem{definition}{Definition}

\numberwithin{equation}{section}

\begin{document}

\DOI{DOI HERE}
\copyrightyear{2021}
\vol{00}
\pubyear{2021}
\access{Advance Access Publication Date: Day Month Year}
\appnotes{Paper}
\copyrightstatement{Published by Oxford University Press on behalf of the Institute of Mathematics and its Applications. All rights reserved.}
\firstpage{1}

%\subtitle{Subject Section}

\title[Time Series]{Time Series Analysis of Traffic Volume }

\author{Xiang Zhang* \footnote{xzhang23@ualberta.ca}, Kexin Fan\footnote{kfan}



\address{\orgdiv{Department of Mathematics and Statistics}, \orgname{University of Alberta}, \country{Canada}}}

\authormark{Xiang et al.}

\corresp[*]{Corresponding author: Xiang Zhang \href{email:xzhang23@ualberta.ca}{xzhang23@ualberta.ca}}


\revised{4}{4}{2022}

%\editor{Associate Editor: Name}

\abstract{This article analyze the .}
\keywords{keyword1;Time Series Analysis; Traffic Data; Keyword4.}

% \boxedtext{
% \begin{itemize}
% \item Key boxed text here.
% \item Key boxed text here.
% \item Key boxed text here.
% \end{itemize}}

\maketitle


\section{Introduction}
The introduction introduces the context and summarizes the manuscript. It is importantly to clearly state the contributions
of this piece of work. 

Nam dui ligula, fringilla a, euismod sodales, sollicitudin\footnote{\url{https://data.gov.uk/}} vel, wisi. Morbi auctor lorem non justo. Nam lacus libero,
pretium at, lobortis vitae, ultricies et, tellus. Donec aliquet, tortor sed accumsan bibendum, erat ligula aliquet magna,
vitae ornare odio metus a mi. Morbi ac orci et nisl hendrerit mollis. Suspendisse ut massa. Cras nec ante. Pellentesque\footnote{\url{https://www.ukdataservice.ac.uk/}}
a nulla. Cum sociis natoque penatibus et magnis dis parturient montes, nascetur ridiculus mus. Aliquam tincidunt
urna. Nulla ullamcorper vestibulum turpis. Pellentesque cursus luctus mauris.

Nam dui ligula, fringilla a, euismod sodales, sollicitudin vel, wisi. Morbi auctor lorem non justo. Nam lacus libero,
pretium at, lobortis vitae, ultricies et, tellus. Donec aliquet, tortor sed accumsan bibendum, erat\footnote{That is not to say that there have been no breaches---researchers using data for a project they are not licensed for or passing data
to unlicensed colleagues---however, this, as far as we are aware of, has never caused harm to individuals or entities in the data.} ligula aliquet magna,
vitae ornare odio metus a mi. Morbi ac orci et nisl hendrerit mollis. Suspendisse ut massa. Cras nec ante. Pellentesque
a nulla. Cum sociis natoque penatibus et magnis dis parturient montes, nascetur ridiculus mus. Aliquam tincidunt
urna. Nulla ullamcorper vestibulum turpis. Pellentesque cursus luctus mauris.



\section{Data}\label{sec2}

We conducted our experiments on Metro Interstate Traffic Volume data set, provided by UC Irvine Machine Learning Repository. 

The traffic data set \cite{pyrkov2018quantitative} is constructed based on hourly interstate 94 Westbound traffic volume data for MN DoT ATR station 301, roughly midway between Minneapolis and St Paul, MN. Recorded data span from October 2012 to October 2018. 


\subsection{Data description}\label{subsec1}

There are nine features for each data entry in provided data set, holiday, temp, rain\_1h, snow\_1h, clouds\_all, weather\_main, weather\_description, data time and traffic volume. Categorical feature holiday is whether it is US National holidays plus regional holiday, Minnesota State Fair on specified date. numerical feature temp is Average temp in kelvin at given time. rain\_1h and snow\_1h measures amount in mm of rain/snow occur in one hour at given time. clouds\_all is numeric percentage of cloud cover. weather\_main is a piece of Short textual description of the current weather. date\_time is day time Hour of the data collected in local CST time. numerical value traffic\_volume reports Hourly I-94 ATR 301 westbound traffic volume. In this experiment, we mostly use feature data\_time and traffic volume to do time series analysis. However, all other features will be exploited to some extend for model fitting and future prediction.  
\subsection{Data preprocessing}\label{subsec2}
Traffic data downloaded from the website is in csv format, with 48205 rows including headings. During early experiments, we found the hourly data is not complete nor consistent. Most recorded days have data for twenty four hours but some comes with less than twenty four hours' data in a day, which results in inconvenience while analyzing the time series relationship.

To tackle the issues, we process the file in two ways. To analyze the monthly/annual trend of traffic volume, we combine different hours across the day and take the average of them, changing data frequency from hourly to daily. This way, fluctuations across a day can be reduced, much like moving average. To analyze the seasonal changing of a day, we selectively pick the time at 6 a.m, 9 a.m, 12 p.m, 3 p.m, 6 p.m and 11 p.m for each day, if data at that time is missing, we use the closest time frame. This way, we can see the change in different time stamp in a day without worrying too much about too frequent changing and dealing with incomplete data. Additionally, we convert all categorical data into numerical values (e.g. 1, 2, 3 \dots )

The data pre-processing is done in Python, and source code for this project is open for public.\footnote{github link}

\subsection{Data visualization}\label{subsec4}



\paragraph{This is an example for fourth level head - paragraph head}

Nam dui ligula, fringilla a, euismod sodales, sollicitudin vel, wisi. Morbi auctor lorem non justo. Nam lacus libero,
pretium at, lobortis vitae, ultricies et, tellus. Donec aliquet, tortor sed accumsan bibendum, erat ligula aliquet magna,
vitae ornare odio metus a mi. Morbi ac orci et nisl hendrerit mollis. Suspendisse ut massa. Cras nec ante.

sd

\section{This is an example for first level head}\label{sec3}

\subsection{This is an example for second level head - subsection head}\label{subsec2}

\subsubsection{This is an example for third level head - subsubsection head}\label{subsubsec2}

Nam dui ligula, fringilla a, euismod sodales, sollicitudin vel, wisi. Morbi auctor lorem non justo. Nam lacus libero,
pretium at, lobortis vitae, ultricies et, tellus. Donec aliquet, tortor sed accumsan bibendum, erat ligula aliquet magna,
vitae ornare odio metus a mi.

\paragraph{This is an example for fourth level head - paragraph head}

 Morbi ac orci et nisl hendrerit mollis. Suspendisse ut massa. Cras nec ante. Pellentesque
a nulla. Cum sociis natoque penatibus et magnis dis parturient montes, nascetur ridiculus mus. Aliquam tincidunt
urna. Nulla ullamcorper vestibulum turpis. Pellentesque cursus luctus mauris.



\section{Equations}\label{sec4}

Equations in \LaTeX{} can either be inline or on-a-line by itself. For
inline equations use the \verb+$...$+ commands. Eg: The equation
$H\psi = E \psi$ is written via the command \verb+$H \psi = E \psi$+.

For on-a-line by itself equations (with auto generated equation numbers)
one can use the equation environment:
\begin{equation}
\|\tilde{X}(k)\|^2 \leq\frac{\sum\limits_{i=1}^{p}\left\|\tilde{Y}_i(k)\right\|^2+\sum\limits_{j=1}^{q}\left\|\tilde{Z}_j(k)\right\|^2 }{p+q}.\label{eq1}
\end{equation}
where,
\begin{align}
D_\mu &=  \partial_\mu - ig \frac{\lambda^a}{2} A^a_\mu, \nonumber \\
F^a_{\mu\nu} &= \partial_\mu A^a_\nu - \partial_\nu A^a_\mu + g f^{abc} A^b_\mu A^a_\nu.\label{eq2}
\end{align}
Notice the use of \verb+\nonumber+ in the align environment at the end
of each line, except the last, so as not to produce equation numbers on
lines where no equation numbers are required. The \verb+\label{}+ command
should only be used at the last line of an align environment where
\verb+\nonumber+ is not used.
\begin{equation}
Y_\infty = \left( \frac{m}{\textrm{GeV}} \right)^{-3}
    \left[ 1 + \frac{3 \ln(m/\textrm{GeV})}{15}
    + \frac{\ln(c_2/5)}{15} \right]
\end{equation}
The class file also supports the use of \verb+\mathbb{}+, \verb+\mathscr{}+ and
\verb+\mathcal{}+ commands. As such \verb+\mathbb{R}+, \verb+\mathscr{R}+
and \verb+\mathcal{R}+ produces $\mathbb{R}$, $\mathscr{R}$ and $\mathcal{R}$
respectively (refer Subsubsection~\ref{subsubsec3}).



\section{Tables}\label{sec5}

Tables can be inserted via the normal table and tabular environment. To put
footnotes inside tables one has to 12pt use the additional ``tablenotes" environment
enclosing the tabular environment. The footnote appears just below the table
itself (refer Tables~\ref{tab1} and \ref{tab2}).


\begin{verbatim}
\begin{table}[t]
\caption{<table-caption>\label{<table-label>}}%
\begin{tabular}{@{}llll@{}}
\toprule
column 1 & column 2 & column 3 & column 4\\
\midrule
row 1 & data 1 & data 2          & data 3 \\
row 2 & data 4 & data 5$^{1}$ & data 6 \\
row 3 & data 7 & data 8      & data 9$^{2}$\\
\botrule
\end{tabular}
\begin{tablenotes}%
\item Source: Example for source.
\item[$^{1}$] Example for a 1st table footnote.
\item[$^{2}$] Example for a 2nd table footnote.
\end{tablenotes}
\end{table}
\end{verbatim}


Lengthy tables which doesn't fit in textwidth should be set as rotated table. For this, we need to use \verb+\begin{sidewaystable}...+ \verb+\end{sidewaystable}+ instead of\break \verb+\begin{table}...+ \verb+\end{table}+ environment.


\begin{table}[!t]
\caption{Caption text\label{tab1}}%
\begin{tabular*}{\columnwidth}{@{\extracolsep\fill}llll@{\extracolsep\fill}}
\toprule
column 1 & column 2  & column 3 & column 4\\
\midrule
row 1    & data 1   & data 2  & data 3  \\
row 2    & data 4   & data 5$^{1}$  & data 6  \\
row 3    & data 7   & data 8  & data 9$^{2}$  \\
\botrule
\end{tabular*}
\begin{tablenotes}%
\item Source: This is an example of table footnote this is an example of table footnote this is an example of table footnote this is an example of~table footnote this is an example of table footnote
\item[$^{1}$] Example for a first table footnote.
\item[$^{2}$] Example for a second table footnote.
\end{tablenotes}
\end{table}

\begin{table*}[h]
\caption{Example of a lengthy table which is set to full textwidth.\label{tab2}}
\tabcolsep=0pt%%
\begin{tabular*}{\textwidth}{@{\extracolsep{\fill}}lcccccc@{\extracolsep{\fill}}}
\toprule%
& \multicolumn{3}{@{}c@{}}{Element 1$^{1}$} & \multicolumn{3}{@{}c@{}}{Element 2$^{2}$} \\
\cline{2-4}\cline{5-7}%
Project & Energy & $\sigma_{calc}$ & $\sigma_{expt}$ & Energy & $\sigma_{calc}$ & $\sigma_{expt}$ \\
\midrule
Element 3  & 990 A & 1168 & $1547\pm12$ & 780 A & 1166 & $1239\pm100$\\
Element 4  & 500 A & 961  & $922\pm10$  & 900 A & 1268 & $1092\pm40$\\
\botrule
\end{tabular*}
\begin{tablenotes}%
\item Note: This is an example of table footnote this is an example of table footnote this is an example of table footnote this is an example of~table footnote this is an example of table footnote
\item[$^{1}$] Example for a first table footnote.
\item[$^{2}$] Example for a second table footnote.\vspace*{-12pt}
\end{tablenotes}
\end{table*}

\begin{sidewaystable}%[!p]
\caption{Tables which are too long to fit, should be written using the ``sidewaystable" environment as shown here\label{tab3}}
\tabcolsep=0pt%
\begin{tabular*}{\textwidth}{@{\extracolsep{\fill}}lcccccc@{\extracolsep{\fill}}}
\toprule%
& \multicolumn{3}{@{}c@{}}{Element 1$^{1}$}& \multicolumn{3}{@{}c@{}}{Element$^{2}$} \\
\cline{2-4}\cline{5-7}%
Projectile & Energy     & $\sigma_{calc}$ & $\sigma_{expt}$ & Energy & $\sigma_{calc}$ & $\sigma_{expt}$ \\
\midrule
Element 3 & 990 A & 1168 & $1547\pm12$ & 780 A & 1166 & $1239\pm100$ \\
Element 4 & 500 A & 961  & $922\pm10$  & 900 A & 1268 & $1092\pm40$ \\
Element 3 & 990 A & 1168 & $1547\pm12$ & 780 A & 1166 & $1239\pm100$ \\
Element 4 & 500 A & 961  & $922\pm10$  & 900 A & 1268 & $1092\pm40$ \\
Element 3 & 990 A & 1168 & $1547\pm12$ & 780 A & 1166 & $1239\pm100$ \\
Element 4 & 500 A & 961  & $922\pm10$  & 900 A & 1268 & $1092\pm40$ \\
Element 3 & 990 A & 1168 & $1547\pm12$ & 780 A & 1166 & $1239\pm100$ \\
Element 4 & 500 A & 961  & $922\pm10$  & 900 A & 1268 & $1092\pm40$ \\
Element 3 & 990 A & 1168 & $1547\pm12$ & 780 A & 1166 & $1239\pm100$ \\
Element 4 & 500 A & 961  & $922\pm10$  & 900 A & 1268 & $1092\pm40$ \\
Element 3 & 990 A & 1168 & $1547\pm12$ & 780 A & 1166 & $1239\pm100$ \\
Element 4 & 500 A & 961  & $922\pm10$  & 900 A & 1268 & $1092\pm40$ \\
\botrule
\end{tabular*}
\begin{tablenotes}%
\item Note: This is an example of table footnote this is an example of table footnote this is an example of table footnote this is an example of~table footnote this is an example of table footnote
\item[$^{1}$] This is an example of table footnote
\end{tablenotes}
\end{sidewaystable}


\section{Figures}\label{sec6}

As per the \LaTeX\ standards one has to use eps images for \verb+latex+ compilation and \verb+pdf/jpg/png+ images for
\verb+pdflatex+ compilation. This is one of the major difference between \verb+latex+
and \verb+pdflatex+. The images should be single page documents. The command for inserting images
for \verb+latex+ and \verb+pdflatex+ can be generalized. The package used to insert images in \verb+latex/pdflatex+ is the
graphicx package. Figures can be inserted via the normal figure environment as shown in the below example:


\begin{figure}[!t]%
\centering
\includegraphics{01}
\caption{This is a widefig. This is an example of long caption this is an example of long caption  this is an example of long caption this is an example of long caption}\label{fig1}
\end{figure}

\begin{figure*}[!t]%
\centering
\includegraphics{02}
\caption{This is a widefig. This is an example of long caption this is an example of long caption  this is an example of long caption this is an example of long caption}\label{fig2}
\end{figure*}


\begin{verbatim}
\begin{figure}[!t]
        \centering\includegraphics{<eps-file>}
        \caption{<figure-caption>}
        \label{<figure-label>}
\end{figure}
\end{verbatim}

For sample purpose, we have included the width of images in the optional argument of \verb+\includegraphics+ tag. Please ignore this. Lengthy figures which doesn't fit in textwidth should be set as rotated mode. For rotated figures, we need to use \verb+\begin{sidewaysfigure}+ \verb+...+ \verb+\end{sidewaysfigure}+ instead of \verb+\begin{figure}+ \verb+...+ \verb+\end{figure}+ environment.

\begin{sidewaysfigure}[!p]%
\centering
\includegraphics{01}
\caption{This is an example for sideways figure. This is an example of long caption this is an example of long caption  this is an example of long caption this is an example of long caption}\label{fig3}
\end{sidewaysfigure}



\section{Algorithms, Program codes and Listings}\label{sec7}

Packages \verb+algorithm+, \verb+algorithmicx+ and \verb+algpseudocode+ are used for setting algorithms in latex.
For this, one has to use the below format:


\begin{verbatim}
\begin{algorithm}
\caption{<alg-caption>}\label{<alg-label>}
\begin{algorithmic}[1]
. . .
\end{algorithmic}
\end{algorithm}
\end{verbatim}


We need to refer above listed package documentations for more details before setting \verb+algorithm+ environment.
To set program codes, one has to use \verb+program+, package. We need to use \verb+\begin{program}+ \verb+...+
\verb+\end{program}+ environment to set program codes. A fast exponentiation procedure:

\begin{algorithm}[!t]
\caption{Calculate $y = x^n$}\label{algo1}
\begin{algorithmic}[1]
\Require $n \geq 0 \vee x \neq 0$
\Ensure $y = x^n$
\State $y \Leftarrow 1$
\If{$n < 0$}
        \State $X \Leftarrow 1 / x$
        \State $N \Leftarrow -n$
\Else
        \State $X \Leftarrow x$
        \State $N \Leftarrow n$
\EndIf
\While{$N \neq 0$}
        \If{$N$ is even}
            \State $X \Leftarrow X \times X$
            \State $N \Leftarrow N / 2$
        \Else[$N$ is odd]
            \State $y \Leftarrow y \times X$
            \State $N \Leftarrow N - 1$
        \EndIf
\EndWhile
\end{algorithmic}
\end{algorithm}

Similarly, for \verb+listings+, one has to use \verb+listings+ package. \verb+\begin{lstlisting}+ \verb+...+ \verb+\end{lstlisting}+ environment is used to set environments similar to \verb+verbatim+ environment. Refer \verb+lstlisting+, package documentation for more details on this.


\begin{minipage}{\hsize}%
\lstset{language=Pascal}% Set your language (you can change the language for each code-block optionally)
\begin{lstlisting}[frame=single,framexleftmargin=-1pt,framexrightmargin=-17pt,framesep=12pt,linewidth=0.98\textwidth]
for i:=maxint to 0 do
begin
{ do nothing }
end;
Write('Case insensitive ');
Write('Pascal keywords.');
\end{lstlisting}
\end{minipage}


\section{Cross referencing}\label{sec8}

Environments such as figure, table, equation, align can have a label
declared via the \verb+\label{#label}+ command. For figures and table
environments one should use the \verb+\label{}+ command inside or just
below the \verb+\caption{}+ command.  One can then use the
\verb+\ref{#label}+ command to cross-reference them. As an example, consider
the label declared for Figure \ref{fig1} which is
\verb+\label{fig1}+. To cross-reference it, use the command
\verb+ Figure \ref{fig1}+, for which it comes up as
``Figure~\ref{fig1}".

\subsection{Details on reference citations}\label{subsec3}

With standard numerical .bst files, only numerical citations are possible.
With an author-year .bst file, both numerical and author-year citations are possible.

If author-year citations are selected, \verb+\bibitem+ must have one of the following forms:


{\footnotesize%
\begin{verbatim}
\bibitem[Jones et al.(1990)]{key}...
\bibitem[Jones et al.(1990)Jones,
                Baker, and Williams]{key}...
\bibitem[Jones et al., 1990]{key}...
\bibitem[\protect\citeauthoryear{Jones,
                Baker, and Williams}
                {Jones et al.}{1990}]{key}...
\bibitem[\protect\citeauthoryear{Jones et al.}
                {1990}]{key}...
\bibitem[\protect\astroncite{Jones et al.}
                {1990}]{key}...
\bibitem[\protect\citename{Jones et al., }
                1990]{key}...
\harvarditem[Jones et al.]{Jones, Baker, and
                Williams}{1990}{key}...
\end{verbatim}}


This is either to be made up manually, or to be generated by an
appropriate .bst file with BibTeX. Then,


{%
\begin{verbatim}
                    Author-year mode
                        || Numerical mode
\citet{key} ==>>  Jones et al. (1990)
                        || Jones et al. [21]
\citep{key} ==>> (Jones et al., 1990) || [21]
\end{verbatim}}


\noindent
Multiple citations as normal:


{%
\begin{verbatim}
\citep{key1,key2} ==> (Jones et al., 1990;
                         Smith, 1989)||[21,24]
        or (Jones et al., 1990, 1991)||[21,24]
        or (Jones et al., 1990a,b)   ||[21,24]
\end{verbatim}}


\noindent
\verb+\cite{key}+ is the equivalent of \verb+\citet{key}+ in author-year mode
and  of \verb+\citep{key}+ in numerical mode. Full author lists may be forced with
\verb+\citet*+ or \verb+\citep*+, e.g.


{%
\begin{verbatim}
\citep*{key} ==>> (Jones, Baker, and Mark, 1990)
\end{verbatim}}


\noindent
Optional notes as:


{%
\begin{verbatim}
\citep[chap. 2]{key}     ==>>
        (Jones et al., 1990, chap. 2)
\citep[e.g.,][]{key}     ==>>
        (e.g., Jones et al., 1990)
\citep[see][pg. 34]{key} ==>>
        (see Jones et al., 1990, pg. 34)
\end{verbatim}}


\noindent
(Note: in standard LaTeX, only one note is allowed, after the ref.
Here, one note is like the standard, two make pre- and post-notes.)


{%
\begin{verbatim}
\citealt{key}   ==>> Jones et al. 1990
\citealt*{key}  ==>> Jones, Baker, and
                        Williams 1990
\citealp{key}   ==>> Jones et al., 1990
\citealp*{key}  ==>> Jones, Baker, and
                        Williams, 1990
\end{verbatim}}


\noindent
Additional citation possibilities (both author-year and numerical modes)


{%
\begin{verbatim}
\citeauthor{key}       ==>> Jones et al.
\citeauthor*{key}      ==>> Jones, Baker, and
                                Williams
\citeyear{key}         ==>> 1990
\citeyearpar{key}      ==>> (1990)
\citetext{priv. comm.} ==>> (priv. comm.)
\citenum{key}          ==>> 11 [non-superscripted]
\end{verbatim}}


\noindent
Note: full author lists depends on whether the bib style supports them;
if not, the abbreviated list is printed even when full requested.

\noindent
For names like della Robbia at the start of a sentence, use


{%
\begin{verbatim}
\Citet{dRob98}      ==>> Della Robbia (1998)
\Citep{dRob98}      ==>> (Della Robbia, 1998)
\Citeauthor{dRob98} ==>> Della Robbia
\end{verbatim}}


\noindent
The following is an example for \verb+\cite{...}+: \cite{rahman2019centroidb}. Another example for \verb+\citep{...}+: \citep{bahdanau2014neural,imboden2018cardiorespiratory,motiian2017unified,american2013acsm,murphy2012machine,ji20123d}.
Sample cites here \cite{krizhevsky2012imagenet,horvath2018dna} and \cite{pyrkov2018quantitative}, \cite{wang2018face}, \cite{lecun2015deep,zhang2018fine,ravi2016deep}.


\section{Lists}\label{sec9}

List in \LaTeX{} can be of three types: numbered, bulleted and unnumbered. ``enumerate'' environment produces numbered list,
``itemize'' environment produces bulleted list and ``unlist'' environment produces unnumbered list.
In each environments, new entry is added via the \verb+\item+ command.
\begin{enumerate}[1.]
\item This is the 1st item

\item Enumerate creates numbered lists, itemize creates bulleted lists and
unnumerate creates unnumbered lists.
\begin{enumerate}[(a)]
\item Second level numbered list. Enumerate creates numbered lists, itemize creates bulleted lists and
description creates unnumbered lists.

\item Second level numbered list. Enumerate creates numbered lists, itemize creates bulleted lists and
description creates unnumbered lists.
\begin{enumerate}[(ii)]
\item Third level numbered list. Enumerate creates numbered lists, itemize creates bulleted lists and
description creates unnumbered lists.

\item Third level numbered list. Enumerate creates numbered lists, itemize creates bulleted lists and
description creates unnumbered lists.
\end{enumerate}

\item Second level numbered list. Enumerate creates numbered lists, itemize creates bulleted lists and
description creates unnumbered lists.

\item Second level numbered list. Enumerate creates numbered lists, itemize creates bulleted lists and
description creates unnumbered lists.
\end{enumerate}

\item Enumerate creates numbered lists, itemize creates bulleted lists and
description creates unnumbered lists.

\item Numbered lists continue.
\end{enumerate}
List in \LaTeX{} can be of three types: enumerate, itemize and description.
In each environments, new entry is added via the \verb+\item+ command.
\begin{itemize}
\item First level bulleted list. This is the 1st item

\item First level bulleted list. Itemize creates bulleted lists and description creates unnumbered lists.
\begin{itemize}
\item Second level dashed list. Itemize creates bulleted lists and description creates unnumbered lists.

\item Second level dashed list. Itemize creates bulleted lists and description creates unnumbered lists.

\item Second level dashed list. Itemize creates bulleted lists and description creates unnumbered lists.
\end{itemize}

\item First level bulleted list. Itemize creates bulleted lists and description creates unnumbered lists.

\item First level bulleted list. Bullet lists continue.
\end{itemize}

\noindent
Example for unnumbered list items:

\begin{unlist}
\item Sample unnumberd list text. Sample unnumberd list text. Sample unnumberd list text. Sample unnumberd list text. Sample unnumberd list text.

\item Sample unnumberd list text. Sample unnumberd list text. Sample unnumberd list text.

\item sample unnumberd list text. Sample unnumberd list text. Sample unnumberd list text. Sample unnumberd list text. Sample unnumberd list text. Sample unnumberd list text. Sample unnumberd list text.
\end{unlist}

\section{Examples for theorem like environments}\label{sec10}

For enunciation like environments, we require \verb+amsthm+ package. 

\begin{theorem}[Theorem subhead]\label{thm1}
Example theorem text. Example theorem text. Example theorem text. Example theorem text. Example theorem text.
Example theorem text. Example theorem text. Example theorem text. Example theorem text. Example theorem text.
Example theorem text.
\end{theorem}

Quisque ullamcorper placerat ipsum. Cras nibh. Morbi vel justo vitae lacus tincidunt ultrices. Lorem ipsum dolor sit
amet, consectetuer adipiscing elit. In hac habitasse platea dictumst. Integer tempus convallis augue.

\begin{proposition}
Example proposition text. Example proposition text. Example proposition text. Example proposition text. Example proposition text.
Example proposition text. Example proposition text. Example proposition text. Example proposition text. Example proposition text.
\end{proposition}

Nulla malesuada porttitor diam. Donec felis erat, congue non, volutpat at, tincidunt tristique, libero. Vivamus
viverra fermentum felis. Donec nonummy pellentesque ante.

\begin{example}
Phasellus adipiscing semper elit. Proin fermentum massa
ac quam. Sed diam turpis, molestie vitae, placerat a, molestie nec, leo. Maecenas lacinia. Nam ipsum ligula, eleifend
at, accumsan nec, suscipit a, ipsum. Morbi blandit ligula feugiat magna. Nunc eleifend consequat lorem.
\end{example}

Nulla malesuada porttitor diam. Donec felis erat, congue non, volutpat at, tincidunt tristique, libero. Vivamus
viverra fermentum felis. Donec nonummy pellentesque ante.

\begin{remark}
Phasellus adipiscing semper elit. Proin fermentum massa
ac quam. Sed diam turpis, molestie vitae, placerat a, molestie nec, leo. Maecenas lacinia. Nam ipsum ligula, eleifend
at, accumsan nec, suscipit a, ipsum. Morbi blandit ligula feugiat magna. Nunc eleifend consequat lorem.
\end{remark}

Quisque ullamcorper placerat ipsum. Cras nibh. Morbi vel justo vitae lacus tincidunt ultrices. Lorem ipsum dolor sit
amet, consectetuer adipiscing elit. In hac habitasse platea dictumst.

\begin{definition}[Definition sub head]
Example definition text. Example definition text. Example definition text. Example definition text. Example definition text. Example definition text. Example definition text. Example definition text.
\end{definition}

Apart from the above styles, we have \verb+\begin{proof}+ \verb+...+ \verb+\end{proof}+ environment - with proof head in italic style and body text in roman font with an open square at the end of each proof environment.

\begin{proof}Example for proof text. Example for proof text. Example for proof text. Example for proof text. Example for proof text. Example for proof text. Example for proof text. Example for proof text. Example for proof text. Example for proof text.
\end{proof}

Nam dui ligula, fringilla a, euismod sodales, sollicitudin vel, wisi. Morbi auctor lorem non justo. Nam lacus libero,
pretium at, lobortis vitae, ultricies et, tellus. Donec aliquet, tortor sed accumsan bibendum, erat ligula aliquet magna,
vitae ornare odio metus a mi.

\begin{proof}[Proof of Theorem~{\upshape\ref{thm1}}]
Example for proof text. Example for proof text. Example for proof text. Example for proof text. Example for proof text. Example for proof text. Example for proof text. Example for proof text. Example for proof text. Example for proof text.
\end{proof}

\noindent
For quote environment, one has to use\newline \verb+\begin{quote}...\end{quote}+
\begin{quote}
Quoted text example. Aliquam porttitor quam a lacus. Praesent vel arcu ut tortor cursus volutpat. In vitae pede quis diam bibendum placerat. Fusce elementum
convallis neque. Sed dolor orci, scelerisque ac, dapibus nec, ultricies ut, mi. Duis nec dui quis leo sagittis commodo.
\end{quote}
Donec congue. Maecenas urna mi, suscipit in, placerat ut, vestibulum ut, massa. Fusce ultrices nulla et nisl (refer Figure~\ref{fig3}). Pellentesque habitant morbi tristique senectus et netus et malesuada fames ac turpis egestas. Etiam ligula arcu,
elementum a, venenatis quis, sollicitudin sed, metus. Donec nunc pede, tincidunt in, venenatis vitae, faucibus vel (refer Table~\ref{tab3}).

\section{Conclusion}

Some Conclusions here.

%%%%%%%%%%%%%%

\appendix

\section{Section title of first appendix}\label{sec11}

Nam dui ligula, fringilla a, euismod sodales, sollicitudin vel, wisi. Morbi auctor lorem non justo. Nam lacus libero,
pretium at, lobortis vitae, ultricies et, tellus. Donec aliquet, tortor sed accumsan bibendum, erat ligula aliquet magna,
vitae ornare odio metus a mi. Morbi ac orci et nisl hendrerit mollis. Suspendisse ut massa. Cras nec ante. Pellentesque
a nulla. Cum sociis natoque penatibus et magnis dis parturient montes, nascetur ridiculus mus. Aliquam tincidunt
urna. Nulla ullamcorper vestibulum turpis. Pellentesque cursus luctus mauris.

\subsection{Subsection title of first appendix}\label{subsec4}

Nam dui ligula, fringilla a, euismod sodales, sollicitudin vel, wisi. Morbi auctor lorem non justo. Nam lacus libero,
pretium at, lobortis vitae, ultricies et, tellus. Donec aliquet, tortor sed accumsan bibendum, erat ligula aliquet magna,
vitae ornare odio metus a mi. Morbi ac orci et nisl hendrerit mollis. Suspendisse ut massa. Cras nec ante. Pellentesque
a nulla. Cum sociis natoque penatibus et magnis dis parturient montes, nascetur ridiculus mus. Aliquam tincidunt
urna. Nulla ullamcorper vestibulum turpis. Pellentesque cursus luctus mauris.

\subsubsection{Subsubsection title of first appendix}\label{subsubsec3}

Example for Unnumbered figure:

\begin{figure}[!h]
\centering
\includegraphics{02}
\end{figure}

Fusce mauris. Vestibulum luctus nibh at lectus. Sed bibendum, nulla a faucibus semper, leo velit ultricies tellus, ac
venenatis arcu wisi vel nisl. Vestibulum diam. Aliquam pellentesque, augue quis sagittis posuere, turpis lacus congue
quam, in hendrerit risus eros eget felis.

\section{Section title of second appendix}\label{sec12}%

Fusce mauris. Vestibulum luctus nibh at lectus. Sed bibendum, nulla a faucibus semper, leo velit ultricies tellus, ac
venenatis arcu wisi vel nisl. Vestibulum diam. Aliquam pellentesque, augue quis sagittis posuere, turpis lacus congue
quam, in hendrerit risus eros eget felis. Maecenas eget erat in sapien mattis porttitor. Vestibulum porttitor. Nulla
facilisi. Sed a turpis eu lacus commodo facilisis. Morbi fringilla, wisi in dignissim interdum, justo lectus sagittis dui, et
vehicula libero dui cursus dui. Mauris tempor ligula sed lacus. Duis cursus enim ut augue. Cras ac magna. Cras nulla.
\begin{figure}[!b]
\centering
\includegraphics{02}
\caption{This is an example for appendix figure\label{fig4}}
\end{figure}

\subsection{Subsection title of second appendix}\label{subsec5}

Sed commodo posuere pede. Mauris ut est. Ut quis purus. Sed ac odio. Sed vehicula hendrerit sem. Duis non odio.
Morbi ut dui. Sed accumsan risus eget odio. In hac habitasse platea dictumst. Pellentesque non elit. Fusce sed justo
eu urna porta tincidunt. Mauris felis odio, sollicitudin sed, volutpat a, ornare ac, erat. Morbi quis dolor. Donec
pellentesque, erat ac sagittis semper, nunc dui lobortis purus, quis congue purus metus ultricies tellus. Proin et quam.
Class aptent taciti sociosqu ad litora torquent per conubia nostra, per inceptos hymenaeos. Praesent sapien turpis,
fermentum vel, eleifend faucibus, vehicula eu, lacus.

Sed commodo posuere pede. Mauris ut est. Ut quis purus. Sed ac odio. Sed vehicula hendrerit sem. Duis non odio.
Morbi ut dui. Sed accumsan risus eget odio. In hac habitasse platea dictumst. Pellentesque non elit. Fusce sed justo
eu urna porta tincidunt. Mauris felis odio, sollicitudin sed, volutpat a, ornare ac, erat. Morbi quis dolor. Donec
pellentesque, erat ac sagittis semper, nunc dui lobortis purus, quis congue purus metus ultricies tellus. Proin et quam.
Class aptent taciti sociosqu ad litora torquent per conubia nostra, per inceptos hymenaeos. Praesent sapien turpis,
fermentum vel, eleifend faucibus, vehicula eu, lacus.

\subsubsection{Subsubsection title of second appendix}\label{subsubsec4}

Lorem ipsum dolor sit amet, consectetuer adipiscing elit. Ut purus elit, vestibulum ut, placerat ac, adipiscing vitae,
felis. Curabitur dictum gravida mauris. Nam arcu libero, nonummy eget, consectetuer id, vulputate a, magna. Donec
vehicula augue eu neque.

\begin{table}[t]%
\caption{This is an example of Appendix table showing food requirements of army, navy and airforce\label{tab4}}%
\begin{tabular*}{\textwidth}{@{\extracolsep\fill}lcc@{}}%
\toprule
col1 head & col2 head & col3 head \\
\midrule
col1 text & col2 text & col3 text \\
col1 text & col2 text & col3 text \\
col1 text & col2 text & col3 text \\
\botrule
\end{tabular*}
\end{table}

Example for an equation inside appendix:
\begin{align}
 p &= \frac{\gamma^{2} - (n_{C} -1)H}{(n_{C} - 1) + H - 2\gamma} \label{1eq:hybobo:pfromgH} \\
 \theta &= \frac{(\gamma - H)^{2}(\gamma - n_{C} -1)^{2}}{(n_{C} - 1 + H - 2\gamma)^{2}} \label{2eq:hybobo:tfromgH}\; .
\end{align}

\section{Example of another appendix section}\label{sec13}%

Nam dui ligula, fringilla a, euismod sodales, sollicitudin vel, wisi. Morbi auctor lorem non justo. Nam lacus libero,
pretium at, lobortis vitae, ultricies et, tellus. Donec aliquet, tortor sed accumsan bibendum, erat ligula aliquet magna,
vitae ornare odio metus a mi. Morbi ac orci et nisl hendrerit mollis. Suspendisse ut massa. Cras nec ante. Pellentesque
a nulla. Cum sociis natoque penatibus et magnis dis parturient montes, nascetur ridiculus mus. Aliquam tincidunt
urna. Nulla ullamcorper vestibulum turpis. Pellentesque cursus luctus mauris
\begin{equation}
\mathcal{L} = i \bar{\psi} \gamma^\mu D_\mu \psi
    - \frac{1}{4} F_{\mu\nu}^a F^{a\mu\nu} - m \bar{\psi} \psi
\label{eq26}
\end{equation}

Nulla malesuada porttitor diam. Donec felis erat, congue non, volutpat at, tincidunt tristique, libero. Vivamus
viverra fermentum felis. Donec nonummy pellentesque ante. Phasellus adipiscing semper elit. Proin fermentum massa
ac quam. Sed diam turpis, molestie vitae, placerat a, molestie nec, leo. Maecenas lacinia. Nam ipsum ligula, eleifend
at, accumsan nec, suscipit a, ipsum. Morbi blandit ligula feugiat magna. Nunc eleifend consequat lorem. Sed lacinia
nulla vitae enim. Pellentesque tincidunt purus vel magna. Integer non enim. Praesent euismod nunc eu purus. Donec
bibendum quam in tellus. Nullam cursus pulvinar lectus. Donec et mi. Nam vulputate metus eu enim. Vestibulum
pellentesque felis eu massa.

Nulla malesuada porttitor diam. Donec felis erat, congue non, volutpat at, tincidunt tristique, libero. Vivamus
viverra fermentum felis. Donec nonummy pellentesque ante. Phasellus adipiscing semper elit. Proin fermentum massa
ac quam. Sed diam turpis, molestie vitae, placerat a, molestie nec, leo. Maecenas lacinia. Nam ipsum ligula, eleifend
at, accumsan nec, suscipit a, ipsum. Morbi blandit ligula feugiat magna. Nunc eleifend consequat lorem. Sed lacinia
nulla vitae enim. Pellentesque tincidunt purus vel magna. Integer non enim. Praesent euismod nunc eu purus. Donec
bibendum quam in tellus. Nullam cursus pulvinar lectus. Donec et mi. Nam vulputate metus eu enim. Vestibulum
pellentesque felis eu massa.

Nulla malesuada porttitor diam. Donec felis erat, congue non, volutpat at, tincidunt tristique, libero. Vivamus
viverra fermentum felis. Donec nonummy pellentesque ante. Phasellus adipiscing semper elit. Proin fermentum massa
ac quam. Sed diam turpis, molestie vitae, placerat a, molestie nec, leo. Maecenas lacinia. Nam ipsum ligula, eleifend
at, accumsan nec, suscipit a, ipsum. Morbi blandit ligula feugiat magna. Nunc eleifend consequat lorem. Sed lacinia
nulla vitae enim. Pellentesque tincidunt purus vel magna. Integer non enim. Praesent euismod nunc eu purus. Donec
bibendum quam in tellus. Nullam cursus pulvinar lectus. Donec et mi. Nam vulputate metus eu enim. Vestibulum
pellentesque felis eu massa.

Nulla malesuada porttitor diam. Donec felis erat, congue non, volutpat at, tincidunt tristique, libero. Vivamus
viverra fermentum felis. Donec nonummy pellentesque ante. Phasellus adipiscing semper elit. Proin fermentum massa
ac quam. Sed diam turpis, molestie vitae, placerat a, molestie nec, leo. Maecenas lacinia. Nam ipsum ligula, eleifend
at, accumsan nec, suscipit a, ipsum. Morbi blandit ligula feugiat magna. Nunc eleifend consequat lorem. Sed lacinia
nulla vitae enim. Pellentesque tincidunt purus vel magna. Integer non enim. Praesent euismod nunc eu purus. Donec
bibendum quam in tellus. Nullam cursus pulvinar lectus. Donec et mi. Nam vulputate metus eu enim. Vestibulum
pellentesque felis eu massa.

%% Example for unnumbered table inside appendix
\begin{table}
\caption{}{%
\begin{tabular}{lcc}%
\toprule
col1 head & col2 head & col3 head \\
\midrule
col1 text & col2 text & col3 text \\
col1 text & col2 text & col3 text \\
col1 text & col2 text & col3 text \\
\botrule
\end{tabular}}{}
\end{table}




\section*{Acknowledgments}
The authors thank the anonymous reviewers for their valuable suggestions. This work is supported in part by funds from the National Science Foundation (NSF: \# 1636933 and \# 1920920).

\section*{Funding}
Funding information here.

%\bibliographystyle{plain}
%\bibliography{reference}

\begin{thebibliography}{10}

\bibitem{bahdanau2014neural}
Dzmitry Bahdanau, Kyunghyun Cho, and Yoshua Bengio.
\newblock Neural machine translation by jointly learning to align and
  translate.
\newblock {\em arXiv preprint arXiv:1409.0473}, 2014.

\bibitem{horvath2018dna}
Steve Horvath and Kenneth Raj.
\newblock Dna methylation-based biomarkers and the epigenetic clock theory of
  ageing.
\newblock {\em Nature Reviews Genetics}, 19(6):371, 2018.

\bibitem{imboden2018cardiorespiratory}
Mary~T Imboden, Matthew~P Harber, Mitchell~H Whaley, W~Holmes Finch, Derron~L
  Bishop, and Leonard~A Kaminsky.
\newblock Cardiorespiratory fitness and mortality in healthy men and women.
\newblock {\em Journal of the American College of Cardiology},
  72(19):2283--2292, 2018.

\bibitem{ji20123d}
Shuiwang Ji, Wei Xu, Ming Yang, and Kai Yu.
\newblock 3d convolutional neural networks for human action recognition.
\newblock {\em IEEE Transactions on Pattern Analysis and Machine Intelligence},
  35(1):221--231, 2012.

\bibitem{krizhevsky2012imagenet}
Alex Krizhevsky, Ilya Sutskever, and Geoffrey~E Hinton.
\newblock Imagenet classification with deep convolutional neural networks.
\newblock In {\em Advances in Neural Information Processing Systems}, pages
  1097--1105, 2012.

\bibitem{lecun2015deep}
Yann LeCun, Yoshua Bengio, and Geoffrey Hinton.
\newblock Deep learning.
\newblock {\em Nature}, 521(7553):436, 2015.

\bibitem{motiian2017unified}
Saeid Motiian, Marco Piccirilli, Donald~A Adjeroh, and Gianfranco Doretto.
\newblock Unified deep supervised domain adaptation and generalization.
\newblock In {\em Proceedings of the IEEE International Conference on Computer
  Vision}, pages 5715--5725, 2017.

\bibitem{murphy2012machine}
Kevin~P Murphy.
\newblock {\em Machine learning: A probabilistic perspective}.
\newblock MIT press, 2012.

\bibitem{american2013acsm}
American~College of~Sports~Medicine et~al.
\newblock {\em ACSM's guidelines for exercise testing and prescription}.
\newblock Lippincott Williams \& Wilkins, 2013.

\bibitem{pyrkov2018quantitative}
Timothy~V Pyrkov, Evgeny Getmantsev, Boris Zhurov, Konstantin Avchaciov,
  Mikhail Pyatnitskiy, Leonid Menshikov, Kristina Khodova, Andrei~V Gudkov, and
  Peter~O Fedichev.
\newblock Quantitative characterization of biological age and frailty based on
  locomotor activity records.
\newblock {\em Aging (Albany NY)}, 10(10):2973, 2018.

\bibitem{rahman2019centroidb}
Syed~Ashiqur Rahman and Donald Adjeroh.
\newblock Centroid of age neighborhoods: A generalized approach to estimate
  biological age.
\newblock In {\em 2019 IEEE EMBS International Conference on Biomedical \&
  Health Informatics (BHI)}, pages 1--4. IEEE, 2019.

\bibitem{ravi2016deep}
Daniele Rav{\`\i}, Charence Wong, Fani Deligianni, Melissa Berthelot, Javier
  Andreu-Perez, Benny Lo, and Guang-Zhong Yang.
\newblock Deep learning for health informatics.
\newblock {\em IEEE {J}ournal of {B}iomedical and {H}ealth {I}nformatics},
  21(1):4--21, 2016.

\bibitem{wang2018face}
Zongwei Wang, Xu~Tang, Weixin Luo, and Shenghua Gao.
\newblock Face aging with identity-preserved conditional generative adversarial
  networks.
\newblock In {\em Proceedings of the IEEE Conference on Computer Vision and
  Pattern Recognition}, pages 7939--7947, 2018.

\bibitem{zhang2018fine}
Ke~Zhang, Na~Liu, Xingfang Yuan, Xinyao Guo, Ce~Gao, and Zhenbing Zhao.
\newblock Fine-grained age estimation in the wild with attention {LSTM}
  networks.
\newblock {\em arXiv preprint arXiv:1805.10445}, 2018.

\end{thebibliography}


%USE THE BELOW OPTIONS IN CASE YOU NEED AUTHOR YEAR FORMAT.
%\bibliographystyle{abbrvnat}
%\bibliography{reference}


\end{document}
